\section{De Te Fabula Narratur}

\begin{quotationx}
The citizen had no freedom of religion; either he participated in the religion of the city or he was banished. The hierarchy of family, tribe and city introduced the idea of a wider and wider influence of the gods, but the knowledge of the one god had been lost. Philosophers and initiates in the mystery schools would have known otherwise, yet they also participated in the sacred ceremonies of the city. The esoteric did not preclude the exoteric, a point emphasized by Guenon and Tomberg.

The practices common to the Ancient City recur in different ways in the Medieval period. Priesthood, prayer, sacrifice, communion meals, liturgy, reading the signs of the times, piety, dispassion, the spiritual integrity and unity of the polis, insider vs outsider recognition, spiritual authority, the primacy of spiritual unity over genetic similarity: these are all aspects of Tradition of the past. The Tradition of the future will look similar.
\flright{\textsc{Cologero}}
\end{quotationx}

Here we see a prediction about the spontaneous arising and arrival of a pre-modern religious tradition, necessarily different in different parts of the earth, but ``at one in heart \& mind'' in their necessary rejection of the dry modern rationalism and liberalism which has proved so virulent since the end of WWII. In our next quote, we will see that Rosenstock-Huessy (the great ``underground man'' of German thought) puts his finger on the same issue, although there is more than a ``whiff of Hegelianism'' about his answer that we will wish to reject.

\begin{quotex}
With a conscious economic organization of the whole earth, subconscious tribal organizations are needed to protect man's mind from commercialization and disintegration. The more our shrinking globe demands technical and economic cooperation, the more necessary it will prove to restore the balance by admitting the primitive archetypes of man's nature... before any tribe or group can sacrifice reason to the unreal myth and magic of pre-history, its food and shelter must be guaranteed by the peaceful world-wide organization of production... Nazism is premature, it cannot co-exist with the possibility of war. Frightened by the Proletarian revolution, Nazism is attempting a `classless society', a solution which lies even beyond the Russian society. They are developing the characteristics of the primitive tribes before they can commit themselves to such an adventure... and the professed pacifism of Hitler hinges upon the fact that the Nazis plan to return into the forests like the Germanic tribes. The Jews, who represent the universal history of mankind, stand in their way. Yet it is perhaps only through the Jews that the world may become a playground for tribal primitivism!
\flright{\textsc{Rosenstock-Huessy}\footnote{\url{https://books.google.com/books?id=4-zkAAAAMAAJ\&q=rosenstock+huessy+out+of+revolution\&dq=rosenstock+huessy+out+of+revolution\&hl=en\&ei=1Ve1TouhBIXhsQK3ssDcAw\&sa=X\&oi=book\_result\&ct=result\&resnum=2\&ved=0CDMQ6AEwAQ}},  ``The Future of Revolution'', \emph{Out of Revolution}}
\end{quotex}

Rosenstock-Huessy was not aware of, and therefore could not possibly have foreseen, the coming useful pre-eminence of esoteric and traditional studies, which will (instead) form the basis of a world-wide relative peace between various competing tribalisms. He did not grasp the extent to which Judaism would identify itself, not with the Judaism of a certain refined type of his day, but with the politically secularized elements of the modern global super-state.

\begin{quotex}
Now, every collapse is different, but the particular shape of the collapse is given not just by the conditions that cause it, but by the quality of our response; men will go mad, or men will go rational, and that will seem like a kind of madness. And it is in this regard that the study of the coming collapse is important for this group, for it is only groups like this, I believe, that can form a remnant, and the purpose of a remnant is to conserve true faith and knowledge, and, when the moment comes, to rebuild the world, or more importantly, to rebuild their own neighborhoods. That moment, I believe, is fast upon us. It is the remnant that must lead the way, must exercise that rational madness which will allow us to rebuild---if not the world---then our communities, and these functioning communities must show the world the way to rebuild itself. Here my thesis is very simple: culture has been subordinated to the needs of commerce, a commerce that has exhibited some rather peculiar and even demonic needs. Now, at many times in the past, the merchant has moved culture, and this was not always a bad arrangement. Commerce sought to ennoble itself with culture, an arrangement that was often to their mutual benefit, as many of the monuments of Italy give testimony. The merchant, through his patronage of the arts and the Church, sought to lift up his fellow citizens, ennoble his city, and obtain honor for himself. But what is happening today is something quite different. Although something of the old spirit of patronage remains, in the main the vast engines of culture have been turned from uplifting the citizen to degrading him. Indeed, the whole point of the exercise is to turn each of us from being a citizen into being a a pure consumer; that is, from being a person who takes responsibility for himself, his family, and his community, into being a person whose self-respect is invested only in what he buys, and who is directed only by unregulated and easily manipulated passions. We are told that the economy is regulated by ``self-interest,'' but this is a lie. Indeed, it cannot be so, since self-interest is never something known in advance, but rather something discovered by experience. Who among us has not had the experience of getting exactly what we wanted only to find that it wasn't what we wanted at all? And who has not feared the worst, only to find that it was all for the best? No, self-interest is revealed to us, not known in advance. What we can know---and what advertising appeals to---is our desires. Desire can be converted to self-interest only when guided by the intelligence to good ends and disciplined by the virtues to good means. Intelligence and virtue: these are the enemies of any good marketing program. Modern advertising appeals not to our virtues but to our vices. And it has at its beck and call an incredible and bewildering array of technologies capable of intruding into every corner of our lives and our souls. Marketing has displaced philosophy to become the preeminent integrative science of the modern age.
\flright{\textsc{John Medaille}\footnote{\url{http://www.frontporchrepublic.com/2011/07/will-there-be-zombies/}}, \textit{Front Porch Republic}}
\end{quotex}

This is the same truth, articulated from a paleo-conservative point of view. Medaille here contemplates the necessary (thank God) death of the advertising culture as the formative element of the earth. Instead, local communities based on liturgy (as Cologero indicates) and using low-tech/high-tech means (even if the local community is a large city like Curitiba\footnote{\url{https://www.youtube.com/watch?v=hRD3l3rlMpo}}) will be the ``nodes'' of a future earth organization.

What is dying is the world of Descartes.


\flrightit{Posted on 2011-11-05 by Logres}

\begin{center}* * *\end{center}

\begin{footnotesize}
\begin{sffamily}
\texttt{Logres on 2011-11-05 at 13:59 said:}

I'll just note (as scholarly reference) that the idea behind Medaille's article (that highly efficient systems are actually becoming brittle over time) is supported by a book like ``Faster''

\url{http://en.wikipedia.org/wiki/Peter\_Deunov}

``Tightly coupled'' systems which are designed for perfect efficiency (over time) use up and/or replace the older coping mechanisms which made them practicable and possible (that is to say ``humane'') in the first place. Then, disaster ``out of the blue'' strikes, just when every thing is at peak operating capacity.

\hfill

\texttt{logres on 2011-11-05 at 21:13 said:}

Should have been:

\url{https://www.amazon.com/Faster-Acceleration-Just-About-Everything/dp/067977548X}


\hfill

\texttt{Charlotte Cowell on 2011-11-07 at 11:46 said:}

Interesting... .within the past week I recommended some very hot under the colour commentators on an alternative news site (typically anti-Christian) to read the prophecies of Peter Deunov, very relevant for the trials and tribulations of the present age.

\hfill

\texttt{logres on 2011-11-07 at 17:36 said:}

Can you give me more of a context?

\hfill

\texttt{Charlotte Cowell on 2011-11-08 at 14:32 said:}

oh well I check out quite a few sites like beforenews.com and godlikeproductions, which cover news and `news' from an alternative (not mainstream) perspective but also include a lot of conspiracy theories, stuff about natural disasters and so on, generally very apocalyptic in tone. I read them for a few reasons. Partly as our mainstream media is woefully inadequate, partly because now and again I glean something of genuine interest in amongst the dross, plus the heads up on certain types of current affairs sooner than usual, but also because I like to counterbalance some of the more far-out reports and comments with something a bit more measured. I can also be a bit of a Christian missionary in those places (as I was on Facebook for quite a long while until I ditched it as a very bad job) because you get a lot of terrible anti-Christian and anti-Jewish sentiment being expressed, almost always from people who simply don't know what they're talking about. Anyways, I can't remember precisely what they were talking about but they were referring to false or distorted prophecies and using them to explain the abysmal state of affairs the world is in while simultaneously blaming God for it all. I expected they would find Deunov very interesting as he is so genuinely esoteric and highly pertinent in the present situation, given that we ARE actually sailing through the belt of lies, zone of contradictions, realm of illusions, however you want to call it. Or have just got through it, perhaps is more accurate. I doubted they would have even read something by a genuine master of the Rosicrucian persuasion so hopefully some of them found something to guide them onto a better path therein... .well away from AC's Thelema and all the trouble that's caused in the spiritual world.

\hfill

\texttt{logres on 2011-11-08 at 23:14 said:}

That's interesting. Evola thought Catholic dogma to be a sort of shield against premature exposure to higher worlds. I first became interested in traditional thought through him, because it was the first ``pagan'' articulation (and a well done articulation) of a metaphysical Fact that was impartial. I'm interested (though) in the Celtic missal and Celtic Christianity in general, going back to college, when I did a paper on the Synod of Whitby. Some of my ancestry comes from Rosshire (I think); even an evangelical author like JI Packer noted that Celtic Puritanism was a different ``strain'' from rationalistic Puritan thought. It's slow going, to work through these things. In general, Celtic spirituality seems well worth preserving and investigating further. Any sources you have will be helpful. Thanks.

\hfill

\texttt{Charlotte on 2011-11-09 at 09:45 said:}

Yes, Celtic Christianity is a way that poses least problems for those with pagan leanings -- ancient spirits, shall we say. There are a few places to look. Gareth Knight is the most well known scion of the Avalon tradition of which Dion Fortune was one of the most famous members -- he is also still active and responsive and good on the Druidic and Hyperborean tradtions, as is Fortune... . Knight has done a good job of integrating Qabalah and both he and fortune will greatly appeal to those with a rational rather than highly emotive temperament. John Porter offers good spiritual guidance. Also very interesting for those with a strong interest in grail traditions is Wellesley Tudor Pole.

For more heart centred material John O' Donohue is a wonderful writer, very moving -- anam cara.

Going back further, the Venerable Bede is a good starting point.

Other than that I understand why Evola said what he did about Catholic dogma and I think he may be right, although one must bear in mind that while the shield is present, so is the dark night of the soul, which can last decades. Simpler if one has a monastic or missionary calling.

Those who attempt to bridge the gulf between the ancient and modern worlds have a different sort of problem on top of this to do with knowledge, which is where Knight and Fortune are very helpful -- the rationality and dry humour helps one survive the heartbreak. Or at least to believe one survives!!

\hfill

\texttt{charlotte on 2011-11-11 at 17:30 said:}

A few words from the venerable Bede I posted on my website recently:

``I know you think I speak this in a raving fit, but let me inform you it is not so; for I tell you, that I see this house filled with so much light, that your candle there seems to me to be dark. ''

And when still no one regarded what she said, or returned any answer, she added, '' Let that candle burn as long as you will; but take notice, that it is not my light, for my light will come to me at the dawn of the day.''

\hfill

\texttt{Logres on 2011-11-12 at 15:31 said:}

Where is this quote in Bede, and can you link to your site?

\hfill

\texttt{charlotte on 2011-11-12 at 15:53 said:}

It's from Book IV Chapter VIII

\url{http://alchemical-weddings.com/}



\hfill

\end{sffamily}
\end{footnotesize}

